%----------------------------------------------------------------------------------------
% NOTATIE EN FIGUREN VAN DE CURSUS GETALTHEORIE
% Staat in de preambule en niet in een hoofdstuk: ook de lijsten met oplossingen gebruiken dit.
%----------------------------------------------------------------------------------------
% Bewaarde grafiek, geplaatst door maak-pdf.py: bestand en pagina in figuren-cache/
\newcommand{\grafiekbestand}[2]{\includegraphics[page=#2,trim=3bp 3bp 3bp 3bp]{figuren-cache/#1}}

% Getallenverzamelingen
\providecommand{\N}{\mathbb{N}}
\providecommand{\Z}{\mathbb{Z}}
\providecommand{\Q}{\mathbb{Q}}
\providecommand{\R}{\mathbb{R}}
\newcommand{\Zn}[1]{\mathbb{Z}_{#1}}
% Deelbaarheid: a \deelt b (a deelt b) en a \ndeelt b (a deelt b niet)
\newcommand{\deelt}{\mid}
\newcommand{\ndeelt}{\nmid}
% Grootste gemene deler en kleinste gemeen veelvoud: \ggd(a, b), \kgv(a, b)
\newcommand{\ggd}{\operatorname{ggd}}
\newcommand{\kgv}{\operatorname{kgv}}
% Congruentie: a \equiv b \md{n} geeft a ≡ b (mod n); de rest: a \bmod n
\newcommand{\md}[1]{\ (\mathrm{mod}\ #1)}
\newcommand{\e}{\mathrm{e}}

% Zoals stappen, maar met een bredere kolom voor de berekening: voor lange formules
\newenvironment{stappenbreed}{%
  \par\renewcommand{\arraystretch}{1.25}%
  \noindent\begin{tabular}{@{}>{$\displaystyle}p{9.8cm}<{$}@{\hspace{4mm}}>{\RaggedRight\small\itshape\arraybackslash}p{4.9cm}@{}}%
}{\end{tabular}\par}

% Zoals stappen, met eigen kolombreedtes: \begin{stappenvrij}{6.2cm}{3.8cm}. Voor naast een figuur.
\newenvironment{stappenvrij}[2]{%
  \par\renewcommand{\arraystretch}{1.25}%
  \noindent\begin{tabular}{@{}>{$\displaystyle}p{#1}<{$}@{\hspace{3mm}}>{\RaggedRight\small\itshape\arraybackslash}p{#2}@{}}%
}{\end{tabular}\par}

% Tekst of berekening links, figuur rechts: \naast[breedte rechts]{links}{rechts}
% De breedte rechts komt uit de figuur zelf; het optionele argument wordt niet gebruikt.
\newsavebox{\naastbox}
\newcommand{\naast}[3][]{%
  \par\noindent\sbox{\naastbox}{#3}%
  \begin{minipage}[t]{\dimexpr\linewidth-\wd\naastbox-3mm\relax}\vspace{0pt}#2\end{minipage}\hfill
  \begin{minipage}[t]{\wd\naastbox}\vspace{0pt}\usebox{\naastbox}\end{minipage}\par\addvspace{1mm}}
% Een rij figuren, gecentreerd
\newcommand{\figuren}[1]{\par\addvspace{1mm}{\centering #1\par}\addvspace{1.5mm}}
% Opgave met een grafiek rechts: \begin{metgrafiek}[breedte]{tikz} opgave en deelvragen \end{metgrafiek}
% De solution komt na \end{metgrafiek}, zoals bij metprent.
\newsavebox{\grafiekbox}
\newenvironment{metgrafiek}[2][]{%
  \sbox{\grafiekbox}{#2}%
  \begin{minipage}[t]{\dimexpr\linewidth-\wd\grafiekbox-3mm\relax}%
}{%
  \end{minipage}\hfill
  \begin{minipage}[t]{\wd\grafiekbox}\vspace{-2mm}\usebox{\grafiekbox}\end{minipage}\par}

% Rooster met assen: \rooster{xmin}{xstap}{xmax}{ymin}{ystap}{ymax}, met getallen op de assen (gehele waarden).
% \roosterleeg doet hetzelfde zonder getallen, voor eigen labels.
\newcommand{\roosterleeg}[6]{%
  \draw[black!14,very thin,xstep=#2,ystep=#5] (#1,#4) grid (#3,#6);
  \draw[->,grijs] (#1,0) -- ({#3+.4*#2},0) node[right,text=black] {$x$};
  \draw[->,grijs] (0,#4) -- (0,{#6+.4*#5}) node[above,text=black] {$y$};}
\newcommand{\rooster}[6]{%
  \roosterleeg{#1}{#2}{#3}{#4}{#5}{#6}
  \pgfmathtruncatemacro{\rtweex}{#1+#2}\pgfmathtruncatemacro{\rtweey}{#4+#5}
  \foreach \t in {#1,\rtweex,...,#3}{\ifnum\t=0\else\node[below,inner sep=1.5pt,text=black] at (\t,0) {$\t$};\fi}
  \foreach \t in {#4,\rtweey,...,#6}{\ifnum\t=0\else\node[left,inner sep=1.5pt,text=black] at (0,\t) {$\t$};\fi}}
% In de functievoorschriften is \x de veranderlijke.
\newcommand{\kromme}[4][]{\draw[accent,thick,#1] plot[domain=#2:#3,samples=90,smooth] (\x,{#4});}

% GETALLENAS met gehele getallen: \getallenas{van}{tot}{stap labels}. Eenheid: de x-eenheid van de tikzpicture.
\newcommand{\getallenas}[3]{%
  \draw[->,grijs] ({#1-.4},0) -- ({#2+.6},0);
  \foreach \t in {#1,...,#2}{\draw[grijs] (\t,-.08) -- (\t,.08);}
  \pgfmathtruncatemacro{\gtwee}{#1+#3}%
  \foreach \t in {#1,\gtwee,...,#2}{\node[below,inner sep=2pt,font=\scriptsize,text=black] at (\t,-.08) {$\t$};}}
% Een sprong op de getallenas, als boog: \sprong[opties]{van}{tot}{label}
\newcommand{\sprong}[4][]{\draw[->,accent,thick,#1] (#2,.12) to[bend left=50] node[above,font=\scriptsize,text=black,inner sep=1pt] {#4} (#3,.12);}
% Een vakje met een getal (voor de zeef en tabellen modulo n): \getalvak[kleur]{x}{y}{getal}
\newcommand{\getalvak}[4][white]{%
  \fill[#1] (#2,#3) rectangle ++(1,1);
  \draw[lijngrijs] (#2,#3) rectangle ++(1,1);
  \node[font=\scriptsize] at ({#2+.5},{#3+.5}) {$#4$};}
% Een klok met n getallen 0 .. n-1, 0 bovenaan, met de wijzers van de klok mee: \klokmod[straal]{n}
\newcommand{\klokmod}[2][1]{%
  \draw[accent,thick] (0,0) circle (#1);
  \pgfmathtruncatemacro{\klokmax}{#2-1}%
  \foreach \k in {0,...,\klokmax}{%
    \fill[accent] ({90-360*\k/#2}:#1) circle (1.4pt);
    \node[font=\scriptsize] at ({90-360*\k/#2}:{#1+.3}) {$\k$};}}
